% !TEX root = ../main.tex
\lecture{7}{2026-10-30}{Entwicklungsökonomie}

\section{Entwicklungsökonomie \& Einkommenseffekte}
Wie lässt sich Armut am effektivsten bekämpfen? Einkommensverteilung, Sozialeprogramme etc.
\begin{definition}
    [Conditioned-Cash-Transfers (CCTs)] sind Geldzahlungen an Individuen unter bestimmten Bedingungen zur Verwendung dieses Geldes. 
\end{definition}
\textbf{Unconditioned Cashtransfers} wenn angenommen, die Empfänger wissen am besten wie ihre Bedürfnisse zu befriedigen sind (Nutzenmaximieren). Vorteil in Kontrollaufwand gegenüber CCTs.
\subsection{Randomized Control-Trials}
Empirische Forschungsmethode, in der Subjekte in 2 Gruppen zufällig zugeteilt werden. Die Kontrollgruppe bekommt keine Massnahme/Behandlung während die Treatment-Gruppe behandelt wird. So lassen sich die Effekte der Massnahmen klar trennen von anderen Einflüsse

\subsubsection{Randomizierung}
\begin{itemize}
    \item Wenn nicht könnten gemeinsame bestimmte Eingenschaften der Subjekten die Resultate beeinflussen
    \item Einige Eigenschaften sind beobachtbar andere nicht
    \item Selektion der Subjekte würde zur Verzerrung der Resultate führen
\end{itemize}

\subsection{Studie: Short-term impact of unconditional cash-Transfers to the poor (experimental evidence from Kenya)}
\subsubsection{Fragen}
\begin{itemize}
    \item Wofür werden CTs in Kenya ausgegeben? 
    \begin{itemize}
        \item Theorie: Wie verändert sich Optimale Güterbündel mit dem Einkommenseffekt der Zahlungen
    \end{itemize}
    \item Wie hängen Entscheidungen der Haushalte  mit Grösse und zeitliches Anfallen der Zahlungen zusammen?
    \item Unterschiede weibliche und männliche Empfänger der Zahlungen?
    \item Wie ändert sich subjektive Wohlbefinden der Empfänger?
\end{itemize}

\subsubsection{Forschungsdesign}
\begin{itemize}
    \item Einmalige oder monatliche Zahlungen
    \item Erhebung Haushaltsdaten
    \begin{itemize}
        \item Vermögenswerte
        \item 
    \end{itemize}
    \item Kontroll und Treatment-Gruppen
    \begin{itemize}
        \item Gruppe bestehen aus Gemeinden
        \item Kontrolle: Gewisse bekommen nichts
        \item Treatment
        \begin{itemize}
            \item Innerhalb der Treatment Gemeinde, bekommen gewisse Haushalte nichts (Altruismus, Effekte auf andere)
            \item Haushalte die Zahlungen bekommen in männlich und weiblich geteilt und jeweils in monatliche oder einmalige Zahlungen
            \item in einmalig und monatliche Zahlungen wurden jeweils kleine und grössere Geldmengen verteilt 
        \end{itemize} 
    \end{itemize}
\end{itemize}
\subsubsection{Ergebnisse}
\begin{itemize}
    \item Grosser teil in \textit{non-land} assets investiert (Nutztiere, Reparaturen, usw.)
    \item Konsum von Nahrungsmittel steigt  (Nahrungssicherheit)
    \item Selbst-erworbenes Einkommen gestiegen 
    \item Hinweise, dass Mittel produktiv Eingesetzt wurden und nicht einmalig konsumiert  
    \item Daten
    \begin{itemize}
        \item Ausgaben in Güter langfristiger Nutzungsdauer (+25\%)
        \item Metalldächer (16\% der Haushalte auf 40\%)
        \item Nutztiere (+50\%)
        \item Konsumgüter(+23\%)
        \item Einnahmen Selbständigkeit und Landwirtschaft (+33\%)
        \item Bargeldersparnisse  (+92\% weil vorher sehr niedrig)
        \item Alkohol/Tabak (keine Veränderung)
    \end{itemize}
\end{itemize}

Wie hängen Resultate von Grösse und zeitlichem Anfallen der Zahlungen ab?

\begin{itemize}
    \item Grosse Zahlungen (Einmalig): Steigen non-land und Konsumgüter
    \begin{itemize}
        \item Anstieg unterproportional zur Grösse der Zahlungen deutet auf notwendige Güter hin
    \end{itemize}
    \item Monatliche Zahlungen effektiver fÚr Nahrungssicherheit
    \item Einmalige Zahlungen tendenziell eher investiert oder gespart 
\end{itemize}

Gab es Unterschiede in Zahlungsempfänger Mann oder Frau in Kenya?

\begin{itemize}
    \item Nein, nur statisch marginal
    \item subjektive Empfindung von Ermächtigung der Frauen etwas gestiegen
\end{itemize}

Wirkung auf subjektives Wohlbefinden:
\begin{itemize}
    \item subjektive Lebenszufriedenheit signifikant gesteigert
    \item Depressionssymptome gesunken 
    \item Basically alle offensichtliche Ergebnisse
\end{itemize}

Langfristige Effekte der UCTs (wenig offensichtlich als kurzfristige Folgen) \footnote{Was war mit dem indirekten Effekt innerhalb der Treatment-Gemeinde bei den Mitglieder, die nichts erhielten?}
\begin{itemize}
    \item Langfristig gibt es einen positiven Effekt auf non-land assets
    \item Sonstige Effekt abhängig von Vergleich zwischen Haushalte in denselben Gemeinden
    \item Möglicherweise im selben Dorf negative Effekte auf Haushalte die 
\end{itemize}
